\documentclass[11pt,a4paper]{article}
\usepackage[T1]{fontenc}
\usepackage[utf8]{inputenc}
\usepackage{lmodern}
\usepackage[margin=2.5cm]{geometry}
\usepackage{microtype}
\usepackage{amsmath,amssymb}
\usepackage{booktabs,longtable,array,tabularx}
\usepackage{xcolor}
\usepackage[numbers,sort&compress]{natbib}
\usepackage[hidelinks]{hyperref}
\usepackage{url}
\usepackage{enumitem}
\setlist{nosep,leftmargin=*}
\setlength{\emergencystretch}{3em}

% ---------------------------------------------------------------- reproducibility markers
% \rn{value}{target} prints a number and a marker [k]; \repro{target} prints the marker alone, for a
% sentence whose numbers all come from one target. Every marker becomes a row of the table in
% Appendix A (value, page, target). A target is FILE::FUNCTION: a pytest test under tests/, or a
% function of papers/scheme-claims/paper_numbers.py. tests/test_paper_numbers.py checks that every
% target resolves.
\newcounter{repro}
\makeatletter
\newcommand{\reprorec}[2]{%
  \refstepcounter{repro}%
  \textsuperscript{\hyperlink{repro.\therepro}{\scriptsize[\therepro]}}%
  \edef\@tempa{\noexpand\addtocontents{lor}{\string\reprorow{\therepro}{\noexpand\thepage}%
    {\noexpand\unexpanded{\noexpand\unexpanded{\unexpanded{#1}}}}{\detokenize{#2}}}}\@tempa}
\newcommand{\rn}[2]{#1\reprorec{#1}{#2}}
\newcommand{\repro}[1]{\reprorec{(see text)}{#1}}
\newtoks\reprorows
\newcommand{\printrepro}{%
  \begingroup\catcode`\_=12
  \def\reprorow##1##2##3##4{\global\reprorows\expandafter{\the\reprorows
    \hypertarget{repro.##1}{##1} & ##2 & ##3 & \nolinkurl{##4}\\}}%
  \@starttoc{lor}\endgroup
  {\small\begin{longtable}{@{}r r >{\raggedright\arraybackslash}p{4.6cm} >{\raggedright\arraybackslash}p{8.2cm}@{}}
  \toprule \# & page & number & reproduced by\\ \midrule\endhead
  \the\reprorows \bottomrule\end{longtable}}}
\makeatother
\def\UrlBreaks{\do\/\do\.\do\:\do\_}
\newcommand{\todo}[1]{\textcolor{red!70!black}{\textbf{[TODO: #1]}}}
\newcommand{\leanpending}{\todo{Lean replay: pending merge}}

\newcommand{\code}[1]{\texttt{#1}}
\newcommand{\V}{\textsf{VERIFIED}}
\newcommand{\R}{\textsf{REFUTED}}
\newcommand{\RE}{\textsf{REFUTED\_EMPIRICALLY}}
\newcommand{\NR}{\textsf{NOT\_REFUTED}}
\newcommand{\U}{\textsf{UNKNOWN}}
\newcommand{\IC}{\textsf{INCONSISTENT}}
\newcommand{\NA}{\textsf{NOT\_APPLICABLE}}
\newcommand{\UF}{\textsf{UNFORMALIZED}}
\newcommand{\SUF}{\textsf{SUFFICIENT}}
\newcommand{\NEC}{\textsf{NECESSARY}}
\newcommand{\EQV}{\textsf{EQUIVALENT}}

\title{A certifying checker for claims about numerical schemes}
\author{The verifier-physics team}
\date{Draft, November 2026}

\begin{document}
\maketitle

\begin{abstract}
Papers in computational physics make claims about their numerical schemes: the scheme is
conservative, second order, monotone under a CFL condition, total-variation diminishing, entropy
stable, or it moves shocks at the right speed. Reviewers check these claims by reading. We describe
a checker that takes a scheme's update formula and a claim about it and returns a verdict. A
positive or negative verdict comes with a certificate: an exact identity (a telescoping sum, a
Farkas combination, non-negative Bernstein coefficients) or an exact rational counterexample that
anyone can replay. A small trusted table says how each check bears on each claim (sufficient,
necessary or equivalent, with stated preconditions), so a failed sufficient check never refutes and
a passed necessary check never verifies. On a benchmark of
\rn{23}{paper_numbers.py::bench_counts} sentences written the way papers write them, split into
\rn{44}{paper_numbers.py::bench_counts} formal readings, the checker decides
\rn{34}{paper_numbers.py::bench_counts} correctly, leaves \rn{9}{paper_numbers.py::bench_counts}
undecided and \rn{1}{paper_numbers.py::bench_counts} unscored, and gives
\rn{0}{paper_numbers.py::bench_counts} wrong verdicts. That zero is on \emph{our own} benchmark: we
wrote the sentences, chose their formal readings and fixed their expected truth values ourselves.
This is the main threat to validity, and we discuss it with the other limits up front
(Section~\ref{sec:trust}). The exact checks run in Python with rational arithmetic; replaying their
certificates in the Lean proof assistant is being merged as we write \leanpending.
\end{abstract}

% ====================================================================================================
\section{A worked example}\label{sec:example}

We follow one sentence through the whole pipeline:
\begin{quote}
\emph{``Rusanov's scheme for Burgers' equation is monotone for CFL $\le 1$.''}
\end{quote}

\paragraph{The equation and the scheme.} Burgers' equation is the conservation law
$u_t + f(u)_x = 0$ with flux $f(u) = u^2/2$. A \emph{finite-volume scheme} divides the line into
cells of width $\Delta x$ and advances cell averages $u_i^n$ by time steps $\Delta t$. A scheme in
\emph{conservative form} has the update
\begin{equation}\label{eq:update}
u_i^{n+1} = H(u_{i-1}^n, u_i^n, u_{i+1}^n) = u_i^n - \lambda\bigl(F(u_i^n, u_{i+1}^n) - F(u_{i-1}^n, u_i^n)\bigr),
\qquad \lambda = \Delta t/\Delta x,
\end{equation}
where $F(u_L, u_R)$, the \emph{numerical flux}, approximates $f$ at the interface between two cells.
Rusanov's flux \cite{rusanov1961} (also called local Lax--Friedrichs) is
\[
F(u_L, u_R) = \tfrac12\bigl(f(u_L) + f(u_R)\bigr) - \tfrac12\,\alpha\,(u_R - u_L),\qquad
\alpha = \max(|u_L|, |u_R|).
\]
The \emph{CFL number} is $\lambda \max|f'(u)| = \lambda\max|u|$, the number of cells a signal crosses
in one step. A scheme is \emph{monotone} if $H$ is non-decreasing in each of its arguments: raising
any one input value never lowers the output. Monotone schemes matter because of a theorem of
Crandall and Majda \cite{crandallmajda1980}: a conservative, consistent, monotone scheme converges to
the physically correct (entropy) solution.

\paragraph{The formal reading.} The checker reads a scheme as a JSON \emph{spec}. Here it holds the
PDE flux \code{u\^{}2/2} and the numerical flux as written, with $|x|$ spelled
$\sqrt{x^2}$ and $\max(a,b) = (a + b + |a - b|)/2$. The sentence becomes the \emph{named claim}
\code{monotone\_under\_cfl} with bound $c = 1$: for every $\lambda > 0$ and all stencil states with
$\lambda|u_j| \le 1$, $H$ is non-decreasing in each argument.

\paragraph{The check.} $H$ is not a polynomial, because of the absolute values. The check splits
every $|x|$ into the two closed \emph{sign cells} $x \ge 0$ and $x \le 0$. On each cell $H$ is a
polynomial, and so are its three partial derivatives. In the scaled variables $w_j = \lambda u_j$ the
partials and the cell boundaries are affine, so on each cell the check looks for an exact
\emph{Farkas identity}: the partial written as a non-negative constant plus a non-negative
combination of the inequalities that define the cell and the CFL box. Such an identity proves the
partial is $\ge 0$ on the whole cell.

At $c = 1$ the search fails and the check looks for a counterexample instead. It returns, in exact
rational arithmetic, $\lambda = 1$ and the state $(u_{i-1}, u_i, u_{i+1})$ =
(\rn{$39/40, 19/20, 1/40$}{paper_numbers.py::rusanov_bounds}): raising $u_i$ by
$1/100$ lowers $H$ from \rn{$2433/3200$}{paper_numbers.py::rusanov_bounds} to
\rn{$60481/80000$}{paper_numbers.py::rusanov_bounds}. All states lie in the CFL box, so the claim is
false. Before the check existed we had found a different counterexample by hand, at
$\lambda = 9/10$:
\[
H(1, 1, -\tfrac{9}{10}) = \text{\rn{$\tfrac{751}{4000}$}{tests/test_schemes.py::test_rusanov_monotone_bound_is_one_third}}
\;>\;
H(1, \tfrac{101}{100}, -\tfrac{9}{10}) = \text{\rn{$\tfrac{18011}{100000}$}{tests/test_schemes.py::test_rusanov_monotone_bound_is_one_third}},
\]
with CFL number \rn{$909/1000$}{paper_numbers.py::rusanov_hand_pair}~$\le 1$. A test evaluates both
values in the spec's own normal form and in an independent implementation of $\max$, and they agree.

\paragraph{The sharp bound.} The same check, asked about $c = 1/3$, returns \V. It covers the state
space with \rn{32}{paper_numbers.py::rusanov_bounds} sign cells and certifies all
\rn{96}{paper_numbers.py::rusanov_bounds} partial derivatives (three per cell) by Farkas identities.
On the cell $u_i \ge u_{i+1} \ge 0$, $u_i \ge u_{i-1} \ge 0$, for example, the identity for the
middle partial is\repro{paper_numbers.py::rusanov_cell_example}
\[
\frac{\partial H}{\partial u_i} = 1 - 2w_i + \tfrac12 w_{i-1} + \tfrac12 w_{i+1}
= \tfrac13 + \tfrac12\, w_{i+1} + \tfrac12\, w_{i-1} + 2\bigl(\tfrac13 - w_i\bigr),
\]
where every term on the right is non-negative on the cell and in the box $|w_j| \le 1/3$. Asked about
$c = 1/2$ and $c = 334/1000$, the check returns \R, each time with an exact decreasing pair
(at $c = 334/1000$: $\lambda = 1$, $u_{i\pm1} = -834499/2500000$,
$u_i = 834833/2500000$, raised by $1/100000$).\repro{paper_numbers.py::rusanov_bounds} So the sharp
bound lies in $[1/3, 334/1000)$. We had guessed $1/2$ (Section~\ref{sec:corrected}).

\paragraph{The verdict.} A trusted table (Section~\ref{sec:judge}) says that the monotonicity check is
\EQV{} to the claim \code{monotone\_under\_cfl}. The judge therefore returns \R{} (exact) for the
sentence at $c = 1$ and \V{} (exact) at $c = 1/3$.\repro{paper_numbers.py::rusanov_claims} In the
benchmark this sentence is C07, and it has three readings, because ``Rusanov'' and ``monotone'' are
both ambiguous. The three-point scheme with local $\alpha$ is monotone under CFL $\le 1$: \R, expected
false. The Rusanov \emph{flux} is monotone ($F$ non-decreasing in $u_L$, non-increasing in $u_R$):
\V, expected true, through a stated bridge lemma. A variant with one global
$\alpha \ge \max|u|$ is monotone under $\lambda\alpha \le 1$: expected true, but the spec language
cannot state the side condition, so it is a gap (\UF).\repro{paper_numbers.py::bench_counts}
Because the scheme is also certified conservative and consistent, the table's Crandall--Majda row
turns the monotonicity certificate at $c = 1/3$ into a verified claim that the fully discrete scheme
selects the entropy solution.\repro{tests/test_schemes.py::test_harten_is_not_refuted_by_tadmor}

\paragraph{The Lean replay.} The refutation is a concrete statement: two rational states and a
rational $\lambda$ at which $H$ decreases. In the Lean replay, the judge writes the negation of the
monotonicity statement for the spec's $H$, serialised from the spec node for node, and a proof
evaluates both sides. Theorem name and replay statistics: \leanpending.

\medskip\noindent
That is the whole pipeline: a sentence, a formal reading, an exact check that returns a witness, a
fixed table that turns the check into a verdict, and a replay of the witness by an independent
checker. The rest of the paper describes each step in general.

% ====================================================================================================
\section{Introduction}\label{sec:intro}

Compilers check code, and proof assistants check mathematics. Claims about numerical schemes in
computational physics papers (``our scheme is conservative and second order, and it captures shocks
at the right speed'') are checked by reviewers reading the paper. Many of these claims are statements
about a finite formula: the update rule of the scheme. Conservation is an algebraic identity.
Consistency and truncation order are statements about a Taylor expansion. Monotonicity, positivity and
the total-variation-diminishing (TVD) property are inequalities over the states and the CFL number.
Such statements can be decided, or at least certified one way or the other, by exact computation.

Other claims are about limits as the grid is refined: convergence order, the speed of computed shocks,
selection of the entropy solution. These follow from structural properties through classical theorems
(Lax--Richtmyer \cite{laxrichtmyer1956}, Lax--Wendroff \cite{laxwendroff1960}, Crandall--Majda
\cite{crandallmajda1980}, Harten \cite{harten1983}), or they can only be tested numerically.

We built a checker that returns a verdict for a claim together with a certificate, and that keeps
the two kinds of reasoning apart. Its parts:
\begin{itemize}
\item \textbf{Exact checks} on the update formula, in rational arithmetic, each returning a witness:
  a recovered flux, a Farkas identity, a list of Bernstein coefficients, or an exact counterexample.
\item \textbf{Empirical studies} (grid refinement, shock speed, convergence to a rarefaction fan),
  labelled as such. An empirical pass never verifies a claim.
\item \textbf{A trusted relation table} (\code{physverify/schemes/claims.py}) that says, for each named
  claim, which checks bear on it, whether each is \SUF, \NEC{} or \EQV, which theorem the row rests
  on, and which other checks must be exactly verified first.
\item \textbf{A benchmark} of sentences phrased as papers phrase them, each split into formal
  readings with an expected truth value fixed in advance.
\item \textbf{A Lean replay} of the exact certificates, so that the identities and counterexamples are
  checked by a small kernel rather than by our Python \leanpending.
\end{itemize}
The design follows the idea of certifying algorithms \cite{mcconnell2011}: the search can be
complicated and untrusted, as long as its output carries a witness that a simple checker can confirm.

On our benchmark the checker never certifies a false claim and never refutes a true one: 0 wrong
among \rn{44}{paper_numbers.py::bench_counts} readings. We wrote that benchmark ourselves, which
limits what the number shows (Section~\ref{sec:trust}).

% ====================================================================================================
\section{Limitations and what is trusted}\label{sec:trust}

We state the limits first, because they bound every result that follows.

\paragraph{The trust table.} This is the table from the repository README, copied verbatim. It
covers the whole verifier, of which the scheme checker is one part.

\begin{center}\small
\begin{tabular}{@{}>{\raggedright\arraybackslash}p{5.6cm}>{\raggedright\arraybackslash}p{2.6cm}>{\raggedright\arraybackslash}p{6.2cm}@{}}
\toprule
Piece & Trusted & Where \\ \midrule
Theory and evidence format & yes & \code{physverify/theory.py} \\
Obligations: statement, guarantee, assumptions & yes & \code{physverify/obligations.py} \\
Judge & yes & \code{physverify/judge.py} \\
Lean definitions and theorems & kernel-checked & \code{lean/PhysVerify/} \\
Interval replay checker for trajectories & validated numerics & \code{physverify/tools/enclosure.py} \\
Tools and planner that search for proofs & no & \code{physverify/tools/}, \code{physverify/planner.py} \\
\bottomrule
\end{tabular}
\end{center}
The scheme statement generator and proof planner get their own rows when the Lean replay merges
\leanpending. For the scheme checker itself, the trusted pieces today are the exact checks in
\code{physverify/schemes/} (their verdicts rest on Python \code{Fraction} arithmetic) and the relation
table in \code{claims.py}.

\paragraph{What stays outside the Lean kernel.} Even once the replay lands, these steps are not
checked by Lean:
\begin{itemize}
\item \textbf{The cell decomposition and its cover.} For the cell-based checks (monotonicity, TVD,
  depth positivity) each cell's fact is a separate statement. The argument that the cells cover the
  state space, the pruning of lower-dimensional cells, and the derivation of each cell's polynomial are
  exact Python. So is the step from the cell facts to the claim (a piecewise-polynomial $H$ with
  non-negative partials on closed covering cells is non-decreasing; Harten's lemma).
\item \textbf{Statement generation.} Python writes the Lean statement from the spec (the update, the
  fluxes, the binders, the domain). A bug there could make Lean prove the wrong statement.
\item \textbf{The claim-reading step.} A person chooses the formal reading of an English sentence, the
  relation between the reading and the named claim, and any bridge lemma that connects them. This is
  formalisation, and nothing checks it except review.
\item \textbf{Relaxations.} The exact depth-positivity check for shallow water replaces
  $\lambda\sqrt{g h_j}$ by a free variable $z_j \ge 0$ and forgets $z_j^2 = \lambda^2 g h_j$. A
  weaker hypothesis gives a stronger theorem, so this is sound, but the relaxation is argued on paper.
\item \textbf{The theorems behind the table rows.} Lax--Richtmyer, Lax--Wendroff, Crandall--Majda,
  Harten's lemma and Tadmor's entropy theory are premises carried into each verdict as assumptions.
  None of them is formalised by us.
\item \textbf{Empirical studies} run in float64 on chosen data and grids.
\end{itemize}

\paragraph{The semi-discrete caveat.} Tadmor's entropy check \cite{tadmor1987,tadmor2003} certifies
an inequality for the semi-discrete scheme (the method of lines, with exact time integration). It says
nothing about forward-Euler time stepping. In this paper every entropy result keeps that label:
``exact, semi-discrete only'', and once replayed, ``proven (Lean), semi-discrete only''. The verdict
label written by the current Lean replay code drops this caveat when it upgrades a verdict to
``proven (Lean)''. A fix is planned; until then, read every entropy-stability label as semi-discrete.

\paragraph{The benchmark is ours.} The main threat to validity is that we wrote the benchmark: the
\rn{23}{paper_numbers.py::bench_counts} sentences, their
\rn{44}{paper_numbers.py::bench_counts} readings and the expected truth values. We tried to phrase
the sentences as papers do and fixed each expected value before running the checker on it, with a
citation or an argument, but we also built the checker. Its blind spots and ours may coincide. A
sentence that a checker of this kind would get wrong is, almost by construction, one we would be less
likely to write. In one case, C07, the expected value differs from the one first written in our task
notes, because of the counterexample in Section~\ref{sec:example}. ``0 wrong'' is evidence that the
checks and the table are consistent with our reading of the literature on these 44 readings. It is
not an error rate on claims from published papers. An independent benchmark, written by people who
have not seen the code, is the obvious next step.

\paragraph{Scope.} One space dimension. Scalar conservation laws, plus $2\times2$ systems for
hyperbolicity, conservation, consistency and depth positivity (shallow water). Three-point stencils,
and five-point stencils for limited schemes of linear advection. Polynomials, $|\cdot|$, $\min$,
$\max$, $\mathrm{minmod}$, and square roots of even powers; division only by monomials. Anything else
returns \U, which is never a wrong verdict. Global order of accuracy for nonlinear schemes is
empirical only.

% ====================================================================================================
\section{Claims and readings}\label{sec:readings}

An English sentence about a scheme usually has more than one sensible meaning. The benchmark
(\code{benchmark/scheme-claims/claims.json}) records each sentence with one or more \emph{readings}.
Each reading is a conjunction of \emph{components}, and each component is either
\begin{itemize}
\item a named claim from the trusted table, applied to a spec file (sometimes with an override, such as
  a different CFL range), with a \emph{relation} to the reading: \code{same} (the named claim is the
  conjunct), \code{necessary} (the conjunct implies the named claim, so only a refutation transfers) or
  \code{sufficient} (the named claim implies the conjunct, so only a verification transfers); or
\item an \code{unformalized} note, saying why no named claim expresses it.
\end{itemize}
A component may carry a \emph{bridge}: a stated lemma that links the sentence's wording to the named
claim. Each reading has an expected truth value (true, false, or ``unsettled'', which is not scored).
The reading's verdict is the conjunction of its components' verdicts.

\paragraph{C12: three readings of one sentence.} ``The Lax--Friedrichs scheme for Burgers' equation is
first-order accurate, conservative, and captures shocks at the correct speed.'' The shock clause can
mean two different things:
\begin{enumerate}
\item \emph{All three clauses, unconditionally}: order 1 on smooth solutions, flux form, and the
  scheme's own shocks move at the Rankine--Hugoniot speed. Expected true; verdict \NR, a gap. The order
  clause is empirical only, since no theorem turns a refinement study of a nonlinear scheme into a
  proof.
\item \emph{The shock clause, conditionally}: the scheme is in flux form, so every bounded limit is a
  weak solution and its discontinuities move at the Rankine--Hugoniot speed. This is the Lax--Wendroff
  theorem \cite{laxwendroff1960}; it does not say the scheme converges. Expected true; \V.
\item \emph{The shock clause, unconditionally}: the scheme's shocks move at the Rankine--Hugoniot speed
  as $\Delta x \to 0$. Through Crandall--Majda this follows from exact monotonicity, conservation,
  consistency and a Lipschitz flux, provided the shock study's $\Delta t/\Delta x$ lies in the
  certified monotone range. Expected true; \V.\repro{paper_numbers.py::bench_counts}
\end{enumerate}

\paragraph{C19: four readings.} ``Our upwind scheme is first-order accurate.'' (1) The global error is
$O(\Delta x)$ in $\ell^2$ or $L^1$ under refinement at fixed $\Delta t/\Delta x$: \V, through the
Lax--Richtmyer row. The bridge records the judgement that this is what ``first-order accurate'' usually
means for a linear scheme. (2) Order 1 in $L^1$ also for discontinuous solutions: expected false
(the rate is $O(\sqrt{\Delta x})$ \cite{kuznetsov1976,tangteng1995}), but the order study has only
smooth test problems, so the reading is \UF. (3) The local truncation error is $O(\Delta x)$: \V,
exactly. (4) Order 1 in $\Delta x$ and in $\Delta t$ under separate refinements, with $\Delta t$ small
enough that the CFL condition holds at the finest grid: expected true, but Lax--Richtmyer covers only
fixed $\Delta t/\Delta x$, so only the refinement studies bear on it: \NR, a
gap.\repro{paper_numbers.py::bench_counts}

\paragraph{Formula dependence.} C17 shows that the reading step includes the choice of formula.
Godunov's flux for Burgers, transcribed literally with nested $\max$, gives \U{} for consistency and
entropy stability: the normaliser leaves $\sqrt{u^4/4}$ instead of $u^2/2$. Rewritten by hand as
$\max(a, -b)^2/2$, the same flux is \V. \U{} is not a wrong verdict, but a model formalising claims
automatically would hit this first.\repro{paper_numbers.py::bench_counts}

% ====================================================================================================
\section{The judge}\label{sec:judge}

\paragraph{Checks and guarantees.} A \emph{check} is a function of the spec that returns a status
(\V, \R, \NR, \RE, \U, \NA), a witness and a list of assumptions. Each check has a \emph{guarantee},
fixed in the table and not read from the check's result: \emph{exact} (rational arithmetic, decides
only through an identity or an exact counterexample) or \emph{empirical} (a float64 study on chosen
data). An empirical check can at best report \NR, and its failure is \RE.

\paragraph{The relation table.} For each named claim, \code{claims.py} lists the checks that bear on
it. Each row has a relation and a premise (the theorem the row rests on, carried into the verdict as an
assumption), and optionally a list of other exact checks that must be \V{} on the same spec for the row
to count at all. One check result gives the claim a polarity:
\begin{center}\small
\begin{tabular}{@{}lll@{}}
\toprule
relation & check passes & check fails \\ \midrule
\SUF & verified & unknown \\
\NEC & not refuted & refuted \\
\EQV & verified & refuted \\ \bottomrule
\end{tabular}
\end{center}
A ``verified'' from an empirical check is downgraded to ``not refuted''. \U{} and \NA{} give unknown.
Over all rows for one claim, any refutation refutes and any verification verifies. If both happen the
verdict is \IC. That means a bug in a check, a spec or the table, never a result about the physics: the
command line exits with code 3 and the tests fail.\repro{tests/test_schemes.py::test_contradictory_evidence_is_inconsistent_and_loud}
The verdicts form a small lattice:
\[
\V,\ \R \ \text{(exact)} \qquad \NR,\ \RE \ \text{(empirical)} \qquad \U \qquad \IC.
\]
The benchmark adds \UF{} for readings no named claim expresses.

\paragraph{Why a failed sufficient check decides nothing.} Tadmor's entropy inequality for the
quadratic entropy is sufficient for ``the method-of-lines scheme selects the entropy solution'', not
necessary. Roe's scheme with Harten's entropy fix fails it exactly (at $(u_L, u_R) = (-1, 1)$), yet on
rarefaction data it converges to the fan. Reported as a verdict on the claim, the raw \R{} would have
said that Harten's fix picks the wrong solution. Through the table, the claim is \U, and the fully
discrete entropy claim is \NR{} from the rarefaction study
alone.\repro{tests/test_schemes.py::test_harten_is_not_refuted_by_tadmor}

\paragraph{Preconditions.} The Crandall--Majda row for \code{selects\_entropy\_solution} requires
conservation and consistency to be exactly verified. The quasilinear upwind scheme for Burgers,
$u_i - \lambda u_i (u_i - u_{i-1})$ for $u_i > 0$ and its mirror image, shows why. It is monotone for
$\lambda\max|u| \le 1/3$ (\V), but it is not conservative, and it moves shocks at the wrong speed. The
row gives \U{} with the unmet requirement ``conservation: \R'', and the claim is instead refuted
exactly by a frozen jump (Section~\ref{sec:certs}).\repro{tests/test_schemes.py::test_crandall_majda_needs_conservation_and_consistency}
The Lax--Richtmyer row for \code{converges\_order\_p} requires the von Neumann check and a structural
scope check (linear update with coefficients in $\Delta t/\Delta x$ only, combined refinement at a
certified $\Delta t/\Delta x$, smooth data). For a nonlinear scheme the scope check returns \NA, and the
row cannot fire.\repro{tests/test_schemes.py::test_lax_richtmyer_never_fires_for_nonlinear_schemes}

\paragraph{Selected rows.} The full table has \rn{21}{paper_numbers.py::table_size} named claims
and \rn{37}{paper_numbers.py::table_size} rows. Some of the rows used in this paper:
\begin{center}\small
\begin{tabularx}{\linewidth}{@{}>{\raggedright\arraybackslash}p{3.6cm}>{\raggedright\arraybackslash}X l l@{}}
\toprule
claim & check (premise) & relation & guarantee \\ \midrule
\code{conservative} & telescoping & \EQV & exact \\
 & shock-speed study & \NEC & empirical \\
\code{monotone\_under\_cfl} & monotonicity on sign cells & \EQV & exact \\
\code{selects\_entropy\_solution} & monotonicity (Crandall--Majda), needs conservation, consistency & \SUF & exact \\
 & steady non-entropy step & \NEC & exact \\
 & frozen jump with $f(u_L) \ne f(u_R)$ & \NEC & exact \\
 & rarefaction study & \NEC & empirical \\
\code{converges\_order\_p} & truncation order (Lax--Richtmyer), needs von Neumann, scope & \SUF & exact \\
 & refinement studies & \NEC & empirical \\
\code{limit\_is\_weak\_solution} & conservation (Lax--Wendroff), needs consistency, Lipschitz flux & \SUF & exact \\
\code{correct\_shock\_speed} & frozen Lax shock & \NEC & exact \\
\code{tvd\_under\_cfl} & Harten's incremental form on sign cells & \SUF & exact \\
 & explicit TV increase & \NEC & exact \\
\bottomrule
\end{tabularx}
\end{center}

% ====================================================================================================
\section{Certificates}\label{sec:certs}

Each exact check returns a witness that a separate checker can replay. All of them work in
\code{Fraction} arithmetic on a normal form of the update, in which $|x| = \sqrt{x^2}$ and the
$\max$/$\min$ functions are written through $|\cdot|$, so no new primitive is needed.

\paragraph{Telescoping (conservation).} A three-point update is conservative on periodic grids if
$H(a, b, c) - b = \Phi(b, c) - \Phi(a, b)$ for some $\Phi$; then the total $\sum_i u_i$ is unchanged.
The check recovers $\Phi$ (it is unique up to a constant) and verifies the identity in the normal form,
or gives an exact periodic state whose total changes. For Lax--Wendroff on linear advection the
recovered flux is
$F(u_L, u_R) = \tfrac12 (u_L + u_R) + \tfrac12\lambda(u_L - u_R)$.\repro{paper_numbers.py::certificate_examples}

\paragraph{Bernstein coefficients in the CFL number.} A linear update
$\sum_k c_k(\nu)\, u_{i+k}$ ($\nu = a\Delta t/\Delta x$) is \emph{positive} (a non-negative combination of
the stencil) on $[\nu_0, \nu_1]$ if each $c_k \ge 0$ there. The check expands each $c_k$ in the
Bernstein basis on the interval \cite{bernstein-farouki2012}; non-negative Bernstein coefficients
prove $c_k \ge 0$ on the whole interval. Lax--Friedrichs on $[-1, 1]$ has Bernstein coefficients
$(0, 1)$ for $c_{-1}$ and $(1, 0)$ for $c_{1}$, so it is
positive.\repro{paper_numbers.py::certificate_examples} Claimed on $[0, 11/10]$, it is \R: the
coefficient of $u_{i+1}$, $(1 - \nu)/2$, is \rn{$-1/20$}{paper_numbers.py::certificate_examples} at
$\nu = 11/10$. Von Neumann stability uses the same idea in two variables: $1 - |g|^2$ is a polynomial
in $\nu$ and $z = \cos\theta$, certified by tensor Bernstein coefficients after stripping the factors
$(1 - z)$ and $(1 + z)$. For FTCS, $1 - |g|^2 = \nu^2 z^2 - \nu^2 = -\nu^2(1 - z^2)$, so it is unstable
for every $\nu \ne 0$; the exact witness is $|g|^2$ = \rn{2}{paper_numbers.py::certificate_examples}
at $\nu = 1$, $z = 0$.

\paragraph{Farkas identities on sign cells.} For nonlinear and limited schemes, the update is
piecewise polynomial. The check splits each $|x|$ lazily (only when the normal form needs its sign) and
drops a cell only with a certificate that it lies in a hyperplane. The retained closed cells cover the
state space. On each cell an inequality $p \ge 0$ is certified by an identity
$p = r + \sum_j \mu_j a_j$ with $r, \mu_j \ge 0$ over the cell's defining inequalities
$a_j \ge 0$ (Farkas' lemma \cite{farkas1902}); for TVD the multipliers $\mu_j(\nu)$ are polynomials in
$\nu$, themselves certified by Bernstein coefficients on pieces of the $\nu$-range. Harten's incremental
form \cite{harten1983} is certified this way for the minmod and superbee flux-limited schemes
\cite{sweby1984} on $\nu \in [0, 1]$: \rn{18}{paper_numbers.py::certificate_examples} sign cells for
minmod and \rn{50}{paper_numbers.py::certificate_examples} for superbee, every half-cell certified.

\paragraph{Exact counterexamples with replay.} A refutation is always a finite object in rational
numbers, replayed by evaluating the scheme's own expression:
\begin{itemize}
\item \emph{Lax--Wendroff is not TVD.} In flux-limiter form on the periodic state
  $u = (0, 0, 0, 0, 0, 1)$ at $\nu = 1/4$, one step raises the total variation from
  \rn{2}{paper_numbers.py::certificate_examples} to
  \rn{$33/16$}{paper_numbers.py::certificate_examples}. A test replays it with the textbook ratio form
  of the limiter.\repro{tests/test_limiters.py::test_tv_increase_replays_with_the_textbook_formula}
\item \emph{The frozen jump $1 \mid 0$.} For the quasilinear upwind scheme on Burgers, the four
  windows across the step $1 \mid 0$ satisfy $H(1,1,1) = H(1,1,0) = 1$ and $H(1,0,0) = H(0,0,0) = 0$ as
  identities in $\lambda$. So the computed solution is the same stationary jump on every grid. But
  $f(1) = 1/2 \ne 0 = f(0)$: the Rankine--Hugoniot speed is
  \rn{$1/2$}{paper_numbers.py::certificate_examples}, and $1 \mid 0$ is a Lax shock
  ($f'(1) > f'(0)$, $f'' \ge 0$). The limit is not a weak solution, so the scheme's shock-speed,
  weak-limit and entropy claims are all \R{}
  exactly.\repro{tests/test_fixed_jump.py::test_quasilinear_upwind_burgers_is_refuted_with_the_pair}
\item \emph{Minmod is TVD but not monotone.} At $\nu = 1/2$,
  $H(0, 4, 7, 8, 0)$ is \rn{$23/4$}{paper_numbers.py::certificate_examples} and
  $H(0, 4, 7, 33/4, 0)$ is \rn{$183/32$}{paper_numbers.py::certificate_examples}, which is smaller, with both states
  strictly inside one sign cell. This is the barrier of Godunov \cite{godunov1959} and of Harten, Hyman
  and Lax \cite{hartenhymanlax1976}: a monotone scheme is at most first
  order.\repro{tests/test_mono5.py::test_limited_schemes_are_not_monotone}
\item \emph{Roe's flux without an entropy fix is not entropy stable.} At $(u_L, u_R) = (-1, 1)$,
  Tadmor's quantity $Q$ is \rn{$-2/3$}{paper_numbers.py::certificate_examples}: the interface produces
  entropy. (Exact, semi-discrete only.)\repro{tests/test_schemes.py::test_roe_entropy_counterexample}
\end{itemize}

\paragraph{The Lean replay.} The exact certificates are replayed in Lean 4 \cite{lean4} with Mathlib
\cite{mathlib2020}. The judge writes each statement from the spec; an untrusted planner writes only
the tactics, from the check's witness. General lemmas (a telescoping update conserves the total on
every periodic grid; a polynomial with non-negative Bernstein coefficients is non-negative on its
interval; Farkas combinations) are proved once. Lean never changes a verdict: a failed proof leaves the
exact-Python label. Theorem count, accepted count and time: \leanpending. Which checks are not yet
covered: \leanpending.

% ====================================================================================================
\section{The benchmark}\label{sec:bench}

Table~\ref{tab:bench} lists every reading with its expected truth value and the verdict on current
main. The counts:\repro{paper_numbers.py::bench_counts}
\begin{center}\small
\begin{tabular}{@{}lrrr@{}}
\toprule
 & expected true & expected false & total \\ \midrule
decisive and correct (match) & \rn{17}{paper_numbers.py::bench_counts} (all \V)
  & \rn{17}{paper_numbers.py::bench_counts} (16 \R, 1 \RE) & \rn{34}{paper_numbers.py::bench_counts} \\
gap & \rn{7}{paper_numbers.py::bench_counts} & \rn{2}{paper_numbers.py::bench_counts}
  & \rn{9}{paper_numbers.py::bench_counts} \\
wrong & 0 & 0 & \rn{0}{paper_numbers.py::bench_counts} \\
\bottomrule
\end{tabular}
\end{center}
One more reading (C18, reading 4) is ``unsettled'' and not scored. The 9 gaps by verdict:
\rn{4}{paper_numbers.py::bench_counts} \UF, \rn{1}{paper_numbers.py::bench_counts} \U,
\rn{4}{paper_numbers.py::bench_counts} \NR. A run takes under a minute:
\code{uv run python -m physverify.schemes.bench benchmark/scheme-claims/}. A test fails if any verdict
contradicts the expected truth, or if the record in \code{claims.json} drifts from what the checker
gives.\repro{tests/test_claims_bench.py::test_bench_has_no_wrong_verdict_and_matches_what_is_recorded}

{\small
\setlength{\tabcolsep}{4pt}
\begin{longtable}{@{}l >{\raggedright\arraybackslash}p{7.8cm} c l l@{}}
\caption{Every reading of the benchmark. ``exp.'' is the expected truth value (T, F, or u for
unsettled); the verdict and result are from
\code{python -m physverify.schemes.bench} on current main.}\label{tab:bench}\\
\toprule
& reading & exp. & verdict & result \\ \midrule\endfirsthead
\toprule & reading & exp. & verdict & result \\ \midrule\endhead
C01 & Lax--Wendroff: conservative and order 2 on smooth solutions & T & \V & match \\
    & Lax--Wendroff: local truncation error $O(\Delta x^2)$ & T & \V & match \\
C02 & Lax--Wendroff: TVD for every $\nu \in [0, 1]$ & F & \R & match \\
    & the same, in flux-limiter form & F & \R & match \\
C03 & Lax--Wendroff on Burgers: Tadmor entropy stable (semi-discrete) & F & \R & match \\
    & Lax--Wendroff on Burgers: converges to the entropy solution & F & \R & match \\
C04 & upwind: max-norm stable for every $\nu \ge 0$ & F & \R & match \\
    & upwind: von Neumann stable for every $\nu \ge 0$ & F & \R & match \\
C05 & Lax--Friedrichs: max-norm stable for $\nu \in [-1, 1]$ & T & \V & match \\
    & Lax--Friedrichs: von Neumann stable for $\nu \in [-1, 1]$ & T & \V & match \\
C06 & central flux + forward Euler: consistent, conservative, stable for small $\nu$ & F & \R & match \\
C07 & Rusanov, local $\alpha$: monotone under $\lambda\max|u| \le 1$ & F & \R & match \\
    & Rusanov numerical flux is monotone & T & \V & match \\
    & Rusanov with one global $\alpha$: monotone under $\lambda\alpha \le 1$ & T & \UF & gap \\
C08 & Engquist--Osher: monotone under $\lambda\max|u| \le 1$, entropy limit & T & \V & match \\
C09 & Engquist--Osher flux (integral form): Tadmor entropy stable & T & \V & match \\
    & Engquist--Osher flux: entropy stable for every convex entropy & T & \UF & gap \\
C10 & Roe: fully discrete limit is the entropy solution & F & \R & match \\
    & Roe: method-of-lines limit is the entropy solution & F & \R & match \\
C11 & Roe + Harten fix: converges to the fan on transonic rarefactions & T & \NR & gap \\
C12 & Lax--Friedrichs on Burgers: order 1, flux form, correct shock speed & T & \NR & gap \\
    & shock clause, conditional (every bounded limit is a weak solution) & T & \V & match \\
    & shock clause, unconditional & T & \V & match \\
C13 & quasilinear upwind: shocks at the Rankine--Hugoniot speed & F & \R & match \\
    & quasilinear upwind: every bounded limit is a weak solution & F & \R & match \\
C14 & upwind (update form): consistent, flux form, positive on $[0, 1]$ & T & \V & match \\
C15 & Lax--Friedrichs (update form): non-negative coefficients on $[-1, 1]$ & T & \V & match \\
C16 & Lax--Friedrichs on Burgers: Tadmor entropy stable for every $\Delta t/\Delta x$ & F & \R & match \\
    & fully discrete entropy inequality for every $\Delta t/\Delta x$ & F & \UF & gap \\
C17 & Godunov flux, $\max$ transcribed literally: consistent, entropy stable & T & \U & gap \\
    & the same flux rewritten as $\max(a, -b)^2/2$ & T & \V & match \\
C18 & minmod MUSCL: order 2 in $L^1$ without extrema & T & \NR & gap \\
    & minmod MUSCL: order 2 in $L^\infty$ with extrema & F & \RE & match \\
    & minmod MUSCL: truncation error $O(\Delta x^2)$ at every point, extrema included & F & \R & match \\
    & minmod MUSCL: order 2 in $L^1$ on sine data & u & \NR & unscored \\
C19 & upwind: global order 1 at fixed $\Delta t/\Delta x$ ($\ell^2$ or $L^1$) & T & \V & match \\
    & upwind: order 1 in $L^1$ for discontinuous solutions & F & \UF & gap \\
    & upwind: local truncation error $O(\Delta x)$ & T & \V & match \\
    & upwind: order 1 in $\Delta x$ and in $\Delta t$, separate refinements & T & \NR & gap \\
C20 & Rusanov for $u_t + (u^2)_x = 0$: consistent, entropy stable for $\eta = u^2$ & T & \V & match \\
C21 & minmod: TVD on $[0, 1]$ & T & \V & match \\
    & superbee: TVD on $[0, 1]$ & T & \V & match \\
C22 & superbee: TVD on $[0, 3/2]$ & F & \R & match \\
C23 & minmod flux-limited scheme: monotone for $\nu \in (0, 1]$ & F & \R & match \\
\bottomrule
\end{longtable}}

\paragraph{The gaps, plainly.} From \code{GAPS.md}, ranked by how often we think such claims appear
(our judgement, not a measurement):
\begin{itemize}
\item \emph{Global order of nonlinear and limited schemes} (C12 reading 1, C18 reading 1). Refinement
  studies can refute an order but never prove one, and nonlinear schemes have no Lax--Richtmyer theorem.
\item \emph{Entropy claims outside one quadratic entropy} (C09 reading 2: every convex entropy, which
  for scalar laws is the E-flux condition \cite{osher1984}; C16 reading 2: a fully discrete cell entropy
  inequality).
\item \emph{Parameters with side conditions} (C07 reading 3: a global $\alpha \ge \max|u|$).
\item \emph{Regimes the spec cannot state} (C19 reading 2: discontinuous data; C19 reading 4: separate
  space and time refinement).
\item \emph{Empirical-only entropy selection for non-monotone schemes} (C11: Roe with Harten's fix is
  not monotone, so Crandall--Majda does not apply, and only the rarefaction study bears on it).
\item \emph{Formula dependence} (C17 reading 1).
\end{itemize}
Every gap is \U, \NR{} or \UF, never a wrong verdict. Most expected-true claims that we could not
verify are claims about limits; most expected-false claims are refuted exactly, even when nothing could
verify the true version of the sentence.

% ====================================================================================================
\section{When the checker corrected us}\label{sec:corrected}

We list the places where we were wrong and the checker, or a rule we adopted because of it, showed it.
Each item points to the commit or test that records it.

\paragraph{Rusanov's monotone bound.} We expected C07 (``monotone under CFL $\le 1$'') to be true; it
appeared in our task notes as an example of a true claim. A hand computation produced the
counterexample at $\lambda = 9/10$ in Section~\ref{sec:example}. We then estimated the sharp bound as
$1/2$. The monotonicity check refuted $1/2$ with an exact pair and certified
$1/3$.\repro{tests/test_schemes.py::test_rusanov_monotone_bound_is_one_third} The Rusanov spec
added with the check already claimed the sharp $1/3$ (commit \code{31f2218}); the hand-found pair
became a regression test (commit \code{7decba5}).

\paragraph{A ``broken'' centred scheme that was conservative.} For the first shallow-water specs we
wrote a centred, non-conservative-looking momentum update with terms
$u_i(m_{i+1} - m_{i-1}) + m_i(u_{i+1} - u_{i-1})$, meant as a broken example. The conservation check
verified it: the terms telescope by a discrete product rule. The broken spec now uses
$h_i\bigl((u^2)_{i+1} - (u^2)_{i-1}\bigr)$, which does
not.\repro{tests/test_systems.py::test_centred_product_rule_telescopes} The scalar analogue is easy to
miss too: the centred quasilinear Burgers update $u_i - \tfrac{\lambda}{2}u_i(u_{i+1} - u_{i-1})$ is
conservative, with flux $F(u_L, u_R) = u_L u_R/2$.\repro{paper_numbers.py::product_rule}

\paragraph{Grid-commensurate coincidences.} Twice an empirical result looked exact to round-off for a
reason unrelated to the scheme. In the shock-speed study, Engquist--Osher's error at crossing level
$0.5$ was zero because of the symmetry of its discrete profile, not because of accuracy. The endpoint
speed estimate came out at round-off whenever the measuring window moved the profile a whole number of
cells. We replaced the estimator by a least-squares slope over the late window at three levels, and
adopted a rule: every empirical pass is re-run at a deliberately incommensurate parameter (for
example $\Delta t/\Delta x = 37/100$ next to $2/5$) before anyone calls it
exact.\repro{tests/test_schemes.py::test_late_window_removes_the_start_up_offset} Roe's
fixed-point error of $1/4$ on rarefaction data survives at both $2/5$ and
$37/100$.\repro{tests/test_schemes.py::test_roe_keeps_the_expansion_shock} The rule later
confirmed that minmod's max-norm order of about 1.3 at extrema is not a commensurability
artefact.\repro{tests/test_limiters.py::test_linf_study_is_not_a_commensurability_artefact}

\paragraph{A Lax--Wendroff-theorem row that did not exist.} A work brief we wrote formalised the
conditional shock-speed reading of C12 ``via the existing Lax--Wendroff theorem row''. There was no
such row; \code{GAPS.md} had proposed it and said it had not been added. The row was then added as a
new claim, \code{limit\_is\_weak\_solution}, with consistency and a locally Lipschitz numerical flux as
preconditions. The Lipschitz condition is a hypothesis of the theorem that our original proposal left
implicit (commit \code{f1ed980}).\repro{tests/test_shock_speed_rows.py::test_limit_is_weak_solution_row}

\paragraph{A shock-speed refutation without a shock.} The first version of the frozen-jump row refuted
\code{correct\_shock\_speed} for any frozen pair $u_L \mid u_R$ with $f(u_L) \ne f(u_R)$. If the pair is
an expansion, the exact solution is a rarefaction fan with no shock in it: a frozen expansion shows a
wrong limit, not a shock moving at the wrong speed. The row now counts a pair only if it is a Lax shock
($f'' \ge 0$ on $[u_R, u_L]$ by Bernstein coefficients, and $f'(u_L) > f'(u_R)$). A test scheme that
freezes exactly the expansions is refuted on the weak-limit and entropy claims and not on the shock
speed (commit
\code{ffb5982}).\repro{tests/test_fixed_jump.py::test_frozen_expansion_refutes_the_weak_limit_but_not_the_shock_speed}

\paragraph{A spec that was more broken than we meant.} We wrote Lax--Wendroff with its diffusion term
halved as a test for the order check. The exact von Neumann check showed it is also unstable:
$|g|^2$ is \rn{$5/4$}{paper_numbers.py::certificate_examples} at $\nu = 1$, $\cos\theta = 0$. And FTCS
at $\Delta t/\Delta x = 1/2$ to $T = 1/4$ shows order 1 in the refinement study (\NR); only the exact
check sees the instability.\repro{tests/test_schemes.py::test_ftcs_study_can_miss_the_instability}

\paragraph{A branch tested before it was rebased.} An order-of-accuracy branch passed the full test
suite in its own worktree and was merged. The branch predated the benchmark, so the test that compares
the benchmark's recorded verdicts with the checker's was not on it. After the merge, main failed that
test: C01 had become decidable and was still recorded as a gap. The fix re-recorded the benchmark
(commit \code{0657f8d}), and since then the merge gate runs the suite on the tree rebased onto current
main. No verdict was wrong; the record was stale.

% ====================================================================================================
\section{Related work}\label{sec:related}

\paragraph{Formal proofs about specific schemes and general theorems.} Boldo et al.\
\cite{boldo2010} proved in Coq the convergence, with its order, of the centred three-point scheme for
the one-dimensional wave equation under a CFL condition. It is a machine-checked proof for one scheme,
not a tool that takes an arbitrary update formula. Tekriwal, Duraisamy and Jeannin \cite{tekriwal2021}
proved the Lax equivalence theorem in Coq for a general family of finite-difference schemes. It
derives convergence from consistency and stability, but does not decide either property for a given
scheme. Our relation table uses such theorems as premises; a formal library of them would let the
premises themselves be checked.

\paragraph{Symbolic analysis of schemes.} Levandovskyy and Martin \cite{martin2012} generate finite
difference schemes with Gr\"obner bases and derive von Neumann stability conditions with cylindrical
algebraic decomposition, for linear constant-coefficient PDEs. As far as we could tell from the
description, the output is a symbolic condition, not an independently checkable certificate.

\paragraph{Certified polynomial inequalities.} Solovyev and Hales \cite{solovyev2013} verify nonlinear
inequalities in HOL Light with Taylor interval approximations, and discuss Bernstein-basis bounds.
Harrison \cite{harrison2007} verifies nonlinear real formulas in a prover through sums-of-squares
certificates. These are ingredients of the kind we use (our certificates are simpler: affine Farkas
identities and univariate or bivariate Bernstein coefficients). We did not find them applied to the
properties of numerical schemes. dReal \cite{dreal2013} is an example of a tool whose answers are
sound in one direction only ($\delta$-sat does not prove satisfiability); our relation table tracks the
same kind of polarity for each check.

\paragraph{Certifying algorithms.} McConnell et al.\ \cite{mcconnell2011} describe programs that
return a witness with each output, together with a simple checker for the witness. Our exact checks
follow this pattern; the Lean replay plays the role of the checker.

\paragraph{Benchmarks.} CFDLLMBench \cite{cfdllmbench2025} evaluates language models on CFD knowledge,
numerical reasoning and solver workflows, scored by running code and comparing solutions. It does not
ask for verdicts on claims about scheme properties. QuanTemp \cite{quantemp2024} is a benchmark of
real-world numerical claims with true/false/conflicting labels, without PDE discretisations or
formal certificates.

\paragraph{What we did not find.} In a search of the literature (summarised in our research notes),
we found no general checker that takes an arbitrary update formula and returns checkable certificates
for conservation, order, stability, monotonicity, TVD or entropy properties, and no benchmark of
true/false claims about numerical-method properties. A search can miss things.
\todo{related work: a second, independent literature search, especially for computer-algebra tools in
numerical analysis (e.g.\ automated modified-equation and stability analysis) that we have not
reviewed}

\paragraph{The classical theory.} The table rows rest on the Lax--Richtmyer equivalence theorem
\cite{laxrichtmyer1956}, the Lax--Wendroff theorem \cite{laxwendroff1960}, Crandall and Majda's
convergence of monotone schemes \cite{crandallmajda1980}, Harten's TVD lemma \cite{harten1983},
Tadmor's entropy stability \cite{tadmor1987,tadmor2003}, Godunov's order barrier \cite{godunov1959}
and its nonlinear form by Harten, Hyman and Lax \cite{hartenhymanlax1976}. The schemes are from
\cite{rusanov1961,engquistosher1981,roe1981,sweby1984}, and the textbook treatments we follow are
LeVeque \cite{leveque2002} and Toro \cite{toro2009}. Hou and LeFloch \cite{houlefloch1994} analyse why
non-conservative schemes converge to wrong solutions, which our shock-speed study shows numerically
for the quasilinear upwind scheme.

% ====================================================================================================
\section{Outlook}\label{sec:outlook}

\paragraph{Systems beyond $2\times2$.} Hyperbolicity is decided by the discriminant for $2\times2$
systems only, and entropy stability is not ported to systems (it needs vector entropy variables and an
entropy-conservative flux defined by a line integral). Invariant-domain properties (Riemann-invariant
bounds) are the natural next certificates.

\paragraph{More than one dimension.} Conservation and consistency carry over directly. For limited
schemes the number of sign cells grows quickly with the stencil, and the claims papers make are
different (positivity, maximum principles, invariant domains), so new checks are needed.

\paragraph{Reward signals.} A checker that never certifies a false claim can grade a model that writes
schemes and claims about them. A verdict of \U{} or \NR{} must earn nothing, or a model learns to write
claims nothing can decide. The benchmark needs adversarial and independently written entries before it
can support that use.

% ====================================================================================================
\appendix

\section{Reproducing the numbers}\label{app:repro}

Every number in the paper carries a marker [$k$]. The table lists each marker, the page it is on, the
number, and the test or function that reproduces it. A target
\code{tests/}\textit{file}\code{.py::}\textit{test} is a pytest test, run with
\code{uv run pytest -q 'tests/}\textit{file}\code{.py::}\textit{test}\code{'}. A target
\code{paper\_numbers.py::}\textit{name} is a function of \code{papers/scheme-claims/paper\_numbers.py},
run with \code{uv run python papers/scheme-claims/paper\_numbers.py} \textit{name}; it prints what the
checker returns.
``(see text)'' marks a sentence whose numbers all come from one target. All numbers were taken from
current main (commit \code{edf7254}). The merge gate is
\code{uv run pytest -q tests/ --ignore=tests/test\_lean.py}.

\printrepro

\section{Glossary}\label{app:glossary}

\begin{description}[style=nextline]
\item[Scheme, update, stencil] A rule that advances cell values $u_i^n$ one time step:
  $u_i^{n+1} = H(\ldots)$ of the values in a few neighbouring cells (the stencil; here three or five).
\item[$\lambda$, $\nu$, CFL number] $\lambda = \Delta t/\Delta x$; for linear advection with speed $a$,
  $\nu = a\lambda$ (the Courant number). The CFL number $\lambda\max|f'(u)|$ is how many cells a signal
  crosses in one step. A CFL condition bounds it.
\item[Conservative form, numerical flux] An update of the form
  \eqref{eq:update}, with a numerical flux $F$ at each interface. On a periodic grid the total of $u$ is
  then unchanged.
\item[Consistent] $F(u, u) = f(u)$: the numerical flux reduces to the PDE flux on constant data.
\item[Truncation error, order] The error of the update applied to an exact smooth solution; order $p$
  if it is $O(\Delta x^p)$ at fixed $\Delta t/\Delta x$. The \emph{global} order is about the computed
  solution's error, which is a different statement.
\item[Von Neumann stability] For a linear scheme, every Fourier mode's amplification factor $g(\theta)$
  satisfies $|g| \le 1$.
\item[Positive] A linear update whose coefficients are all non-negative.
\item[Monotone] $H$ is non-decreasing in each argument.
\item[TVD] Total variation diminishing: $\sum_i |u_{i+1} - u_i|$ does not increase in a step.
\item[Weak solution, Rankine--Hugoniot] A solution in the integral sense, which may have jumps; a jump
  from $u_L$ to $u_R$ moves at $s = (f(u_L) - f(u_R))/(u_L - u_R)$.
\item[Entropy solution, Lax shock] The physically relevant weak solution, singled out by entropy
  inequalities. For convex $f$ a jump is admissible (a Lax shock) if $f'(u_L) > f'(u_R)$; otherwise the
  entropy solution is a rarefaction fan.
\item[Entropy stable (Tadmor)] The semi-discrete scheme satisfies a discrete entropy inequality for a
  given entropy pair; checked as $Q = -[v](F - F^*) \ge 0$, where $F^*$ is the entropy-conservative flux.
\item[Semi-discrete, fully discrete] With exact time integration (method of lines), or with the actual
  time step.
\item[Riemann problem] Initial data with one jump, $u_L$ for $x < x_0$ and $u_R$ for $x > x_0$.
\item[Spec] The JSON file that holds a PDE, a scheme and the claims made about it.
\item[Named claim] A precise property in the trusted table, such as \code{monotone\_under\_cfl}.
\item[Reading] One formal meaning of an English sentence: a conjunction of named claims on specs, with
  relations and bridges.
\item[Bridge] A stated lemma linking a sentence's wording to a named claim.
\item[Check, witness] A function of the spec that returns a status and the data that justify it.
\item[Certificate] A witness that a simpler checker can verify: an identity, a list of non-negative
  coefficients, or an exact counterexample.
\item[Sign cell] A region where every $|x|$ in the update has a fixed sign, so the update is a polynomial.
\item[Bernstein coefficients] The coefficients of a polynomial in the Bernstein basis of an interval;
  if all are non-negative, the polynomial is non-negative on the interval.
\item[Farkas identity] $p = r + \sum_j \mu_j a_j$ with $r, \mu_j \ge 0$, which proves $p \ge 0$ wherever
  every $a_j \ge 0$.
\item[Guarantee] Exact (rational arithmetic, identity or counterexample) or empirical (float64 study).
\item[SUFFICIENT, NECESSARY, EQUIVALENT] How a check bears on a claim: its pass implies the claim; the
  claim implies its pass; both.
\item[Verdicts] \V, \R{} (exact); \NR, \RE{} (empirical); \U{} (nothing decisive); \IC{} (a bug); and in the
  benchmark \UF{} (no named claim expresses the reading).
\item[Match, gap, wrong] A decisive verdict that agrees with the expected truth; a non-decisive one; a
  decisive one that disagrees.
\end{description}

\bibliographystyle{plainnat}
\bibliography{refs}

\end{document}
